\documentclass[10pt,a4paper]{amsart}
\usepackage[margin=19mm]{geometry}
\usepackage{amsmath,amssymb,mathtools}
\usepackage[scr=rsfs,cal=euler]{mathalfa}
\usepackage{xfrac}
\usepackage[unicode,hidelinks,bookmarks=false]{hyperref}
\newcommand{\ie}{i.e.\ }
\newcommand{\cf}{cf.\ }
\newcommand{\resp}{respectively\ }
\newcommand{\Subsection}{\subsection}
\providecommand{\nobreakdash}{\nobreak\hbox{-}}
\setlength{\emergencystretch}{2em}
\multlinegap0pt
\usepackage{amssymb}
\usepackage{mathtools}
\usepackage[shortlabels]{enumitem}
\usepackage{enumerate}
\usepackage{tikz}
\usepackage{xcolor}
\newtheorem{lemma}{Lemma}
\newcommand{\C}{\mathbb C}
\newcommand{\CP}{\mathbb{CP}}
\title{Semiclassical Limits of Strongly Parabolic Higgs Bundles and Hyperpolygon Spaces}
\author{Lynn Heller, Sebastian Heller, Claudio Meneses}
\date{Source: arXiv:2512.24236v1 (CC BY 4.0)}
\begin{document}
\maketitle

\noindent\emph{Source and licence.} Semiclassical Limits of Strongly Parabolic Higgs Bundles and Hyperpolygon Spaces. Version arXiv:2512.24236v1 (CC BY 4.0), distributed by arXiv under the Creative Commons Attribution 4.0 International licence, http://creativecommons.org/licenses/by/4.0/, as recorded in the arXiv licence metadata for this exact version and shown on the abstract page https://arxiv.org/abs/2512.24236v1. Full licence notice: \texttt{licenses/arxiv-2512.24236-CC-BY-4.0.txt}.\par\smallskip

\noindent\textit{Editorial scope.} Counted unit: the article's lemma lem:sum-P-zero (source0.tex lines 1899--1906). The article's complete original proof of it is delivered with it (source0.tex lines 1908--1967, one \texttt{proof} environment of 60 lines). No mathematics is written, repaired or re-derived: the statement and the proof are the article's own lines, in the article's own notation, and no retained line is substituted or re-flowed.  Retained uncounted as the source-local context the proof needs: lines 1598--1602; lines 1604--1608; lines 1632--1637; lines 1828--1839; lines 1830--1836.  Nothing was omitted from any retained span, so the package carries no omission record.  No retained line contains a citation, so the wrapper carries no bibliography and no bibliography entry.  No retained line contains a figure environment, an image-inclusion command or drawing code, so no image is delivered and the package declares no asset dependency.  Editorial formatting only, in addition: an AMS \texttt{amsart} preamble that restates the article's own declarations and macros used by the retained lines, namely the environments \texttt{lemma} on separate counters, together with the standard packages those lines need.  The retained environments print in this document as Lemma 1 in order of appearance, whereas the article prints the same items with the numbering of its own full document.  The complete native source of the article as distributed in the arXiv e-print for this exact version is bundled unchanged as \texttt{sources/191867/source0.tex}.
\begin{equation}\label{eq:xicentral}
\underline\xi
:=
\sum_{j=1}^n \mathfrak A_j\,\frac{dz}{z-p_j}.
\end{equation}

\begin{equation}\label{eq:defxit}
\xi(t)
=
t\bigl(\underline\xi+\hat\xi(t)\bigr),
\end{equation}

\begin{lemma}\label{lem:IFTt0}
Let $\xi(t)$ be given by \eqref{eq:defxit}, where $\underline\xi$ is defined in
\eqref{eq:xicentral} and $\hat\xi(0)=0$.
Then conditions \emph{(i)}, \emph{(ii)}, and \emph{(iii)} of the monodromy problem
are satisfied at $t=0$.
\end{lemma}

\begin{lemma}\label{lem:modified-monodromy}
The modified monodromy problem \emph{(i)}, \emph{(iia)}, and \emph{(iii)} admits a solution such that
\begin{equation}\label{eq:iftaddon}
\sum_j P_j
\in
\mathfrak{su}(2)\oplus \Lambda^{+}\mathfrak{sl}(2,\C)
\subset
\Lambda^{\geq0}\mathfrak{sl}(2,\C).
\end{equation}
The solution is unique up to (time--dependent) overall $\mathrm{SU}(2)$--conjugation, and can be chosen
to depend real analytically on $t$ and on $\underline\xi$.
\end{lemma}

\begin{equation}
\sum_j P_j
\in
\mathfrak{su}(2)\oplus \Lambda^{+}\mathfrak{sl}(2,\C)
\subset
\Lambda^{\geq0}\mathfrak{sl}(2,\C).
\end{equation}

\begin{lemma}\label{lem:sum-P-zero}
Consider the solution $\hat\xi(t)$ of the modified monodromy problem
in Lemma~\ref{lem:modified-monodromy} such that \eqref{eq:iftaddon} holds.
Then
\[
\sum_j P_j(t)=0.
\]
\end{lemma}

\begin{proof}
Consider the monodromy of $d+\xi(t)$ based at $q$ as a representation
\[
\pi_1(\CP^1\setminus\{p_1,\dots,p_n,\infty\},q)\to\Lambda\mathrm{SL}(2,\C)
\]
of the $(n+1)$--punctured sphere.
Let $\gamma_\infty$ be a simple loop winding once around $z=\infty$ and satisfying
\[
\gamma_\infty * \gamma_n * \dots * \gamma_1 = 1.
\]

We compute the corresponding logarithmic derivative $B_\infty(t)$ in two ways.
On the one hand, it is given by the logarithmic derivative of
\[
M_\infty(t)
=
\bigl(M_n(t)\, M_{n-1}(t)\dots M_1(t)\bigr)^{-1}
\in
\Lambda\mathrm{SL}(2,\C).
\]
Since $B_j(t)\in\Lambda\mathfrak U$ and $M_j(0)=\mathrm{Id}$ for $j=1,\dots,n$,
we have $M_j(t)\in\Lambda\mathcal U$ for $t$ sufficiently small.
Hence $M_\infty(t)\in\Lambda\mathcal U$, and therefore
\[
B_\infty(t)\in\Lambda\mathfrak U.
\]

On the other hand, we expand the residue at $z=\infty$ in terms of $t$ at $t=0$:
\[
\sum_j P_j(t)= \mathcal C^{(k)}\, t^{k}+o(t^{k}),
\]
for some $k\geq1$ and some
\[
\mathcal C^{(k)}\in \mathfrak{su}(2)\oplus \Lambda^{+}\mathfrak{sl}(2,\C)
\]
by Lemma~\ref{lem:modified-monodromy}.
Arguing as in the proof of Lemma~\ref{lem:IFTt0}
(by expanding the parallel frame $\Psi_t$ solving
$d\Psi_t+\xi(t)\Psi_t=0$ with $\Psi_t(q)=\mathrm{Id}$ in $t$),
we obtain (for some $k\geq1$)
\[
B_\infty(t)=2\pi i\, \mathcal C^{(k)}\, t^{k}+o(t^{k}).
\]

Thus
\[
2\pi i\, \mathcal C^{(k)}\in 2\pi i\,
\bigl(\mathfrak{su}(2)\oplus \Lambda^{+}\mathfrak{sl}(2,\C)\bigr)
\cap \Lambda\mathfrak U
=
\{0\},
\]
and since $B_\infty(t)$ is real analytic in $t$, we conclude that
$B_\infty(t)=0$ for all $t$ sufficiently small.
Equivalently,
\[
\sum_j P_j(t)=0.
\]
This completes the proof.
\end{proof}

\end{document}
